\begin{flushleft}
\begin{tikzpicture}[every node/.style={font=\footnotesize},
  chSixRelationMatrix/.style={matrix of nodes,nodes in empty cells,
    nodes={draw=black,line width=.6pt,align=left,inner xsep=3pt,inner ysep=2pt,minimum height=7mm,anchor=center,text height=2.7mm,text depth=.8mm},
    column sep=-.6pt,row sep=-.6pt,row 1/.style={nodes={fill=chSixTint}}}]
\matrix (p) [chSixRelationMatrix,column 1/.style={nodes={text width=23mm}},column 2/.style={nodes={text width=23mm}},column 3/.style={nodes={text width=26mm}}] {
PhoneID (PK) & StaffID (FK) & PhoneNo\\
P01 & S11 & 01700000001\\
P02 & S11 & 01800000002\\
P03 & S12 & 01900000003\\
};
\matrix (s) [chSixRelationMatrix,anchor=north west,column 1/.style={nodes={text width=22mm}},column 2/.style={nodes={text width=15mm}}]
at ([xshift=18mm]p.north east) {
StaffID (PK) & Name\\
S11 & Rima\\
S12 & Fahim\\
};
\node[anchor=south west,font=\small] at ([yshift=2mm]p.north west) {\textbf{Phone: M}};
\node[anchor=south west,font=\small] at ([yshift=2mm]s.north west) {\textbf{Staff: 1}};
\draw[bookarrow] (p-2-3.east)--(s-2-1.west);
\draw[bookarrow] (p-3-3.east)--(s-2-1.west);
\draw[bookarrow] (p-4-3.east)--(s-3-1.west);
\end{tikzpicture}
\captionof{table}{StaffID মিলিয়ে একাধিক Phone Record থেকে একই Staff Record}
\end{flushleft}
